% !TeX root=../main.tex

\chapter{نتیجه گیری}
در این پایان‌نامه، علاوه بر جمع‌آوری و معرفی مجموعه‌داده پزشکی فارسی، دو مدل زبانی کوچک تخصصی پزشکی نیز آموزش داده و ارائه شدند.هر چند تا به کارگیری این مدل های کوچک در محیط واقعی فاصله بسیاری داریم ولی نتایج به‌دست‌آمده نشان‌دهنده عملکرد نسبتا مطلوب این مدل‌ها است و می‌توانند به‌عنوان گام‌های اولیه در مسیر توسعه مدل‌های زبانی کوچک و کارآمد پزشکی به زبان فارسی محسوب شوند.

در نظر داشته باشید که طبق درس تلخ ساتون
\footnote{\lr{sutton's bitter lesson}}
هر اندازه یک مدل یادگیری ماشین بزرگتر باشد و با حجم داده بیشتری تمرین داده شده باشد عملکرد بهتری خواهد داشت بنابراین منطقی است که برای کاربرد های پزشکی که در آن تنها عملکرد مد نظر است بایستی مدل های زبانی با اندازه بسیار بزرگتری تمرین داده شوند که نیاز به سخت افزار بیشتری دارند؛ سخت افزاری که اجاره آن هزاران دلار هزینه در پی خواهد داشت و تهیه آن برای یک دانشجو مقدور نیست.

\section{نوآوری ها}

به طور خلاصه نوآوری های این پایان نامه به شرح زیر است:

\begin{itemize}
    \item 
    معرفی اولین مدل های زبانی پزشکی فارسی منبع‌باز که در مقایسه با سایر مدل‌های قابل اجرا روی دستگاه‌های خانگی (مانند لپ‌تاپ و کامپیوتر شخصی) به نتایج امیدوارکننده دست یافته است.
    \item 
    معرفی یک پیکره پزشکی فارسی که از طریق خزش
    \footnote{\lr{crawling}}
    وب‌سایت‌های مختلف پزشکی جمع‌آوری شده است.
    \item
    معرفی اولین مجموعه‌داده پرسش و پاسخ پزشکی فارسی فرم آزاد
    \footnote{\lr{free form}}
        که از طریق خزش سایت‌های پرسش و پاسخ پزشکی جمع‌آوری شده است.
    
    \item
    ترجمه بخش پزشکی مجموعه داده
    \lr{MMLU}
    به زبان فارسی که می‌تواند به‌عنوان یک معیار استاندارد و قابل اعتماد برای ارزیابی هر مدل زبانی پزشکی فارسی مورد استفاده قرار گیرد.
    \item
    معرفی چارچوب های کارآمد که با استفاده از یک حلقه آموزگار–شاگرد باعث بهبود توانایی دلیل آوری مدل زبانی شاگرد میشوند.
\end{itemize}


\section{
تحقیقات آینده
}

\subsection{
توسعه یک مدل زبانی دیداری
}
در پژوهش‌های انجام‌شده در این پایان‌نامه، تمرکز تنها بر دریافت متن به‌عنوان ورودی بوده است، در حالی که کاربردهای واقعی پزشکی در بسیاری موارد به دریافت تصویر نیز نیاز دارند. اگر پارامترهای به‌روزشده خود را به‌جای مدل آیا اکسپنس
\cite{b71}
، با مدل آیا ویژن
\cite{b74}
که یک مدل زبانی‌ دیداری
\footnote{\lr{vision language model}}
است ادغام کنیم
\footnote{
توجه داشته باشید که مدل‌های معرفی‌شده در این پایان‌نامه با استفاده از
\lr{LoRA}
آموزش داده شده‌اند؛ به این معنا که تنها بخش بسیار کوچکی از پارامترها به‌روزرسانی شده است. این مجموعه‌ی محدود از پارامترهای به‌روزشده، علاوه بر امکان ترکیب با مدل پایه یعنی آیا اکسپنس، قابلیت ادغام در مدل آیا ویژن را نیز دارند.
}
، مدل‌های گائوکرنا وی و گائوکرنا آر نیز قادر خواهند بود ورودی‌های چندوجهی
\footnote{\lr{multimodal}}
از جمله تصاویر پزشکی را دریافت و پردازش کنند. با این حال، توانایی این مدل‌ها در پاسخ‌گویی به پرسش‌های چندوجهی هنوز به‌طور جامع ارزیابی نشده است؛ بنابراین سنجش این قابلیت می‌تواند در اولویت تحقیقات آینده قرار گیرد.


\subsection{
بهبود مدل از طریق یادگیری تقویتی مبتنی بر بازخورد انسانی
}
یکی دیگر از تحقیقات آینده می‌تواند شامل همکاری سازمان‌یافته و مستمر با پزشکان متخصص در حوزه‌های مرتبط باشد تا از تجربه‌های بالینی واقعی آنان برای بهبود عملکرد مدل بهره گرفته شود. این همکاری به‌ویژه در چارچوب یادگیری تقویتی مبتنی بر بازخورد انسانی
\footnote{\lr{reinforcement learning with human feedback}}
اهمیت می‌یابد؛ زیرا بازخورد مستقیم متخصصان می‌تواند مدل را در جهت ارائه پاسخ‌های دقیق‌تر، قابل‌اعتمادتر و سازگار با نیازهای عملی بالین هدایت کند.

\subsection{
توسعه یک مدل مستقل از نوع درخواست
}
همانطور که در فصل ششم اشاره شد مدل گائوکرنا وی در مواجهه با درخواست‌های
\footnote{\lr{prompt}}
مستقیم عملکرد بهتری دارد، در حالی که مدل گائوکرنا آر در در مواجهه با درخواست‌های زنجیره‌استدلال
\footnote{\lr{chain of thought}}
عملکرد برتری ارائه می‌کند. این موضوع بیانگر آن است که هر دو مدل به‌شدت وابسته به نوع پرامپت هستند و پژوهش‌های آینده باید این محدودیت مهم را برطرف کنند.

\subsection{
بررسی عدم اطمینان در مدل آیا اکسپنس
}
همان‌طور که در فصل پنجم بررسی شد، مدل زبانی آیا اکسپنس پیشرفت قابل‌توجهی در حوزه زبان‌های کم‌منبع
\footnote{\lr{low resource languages}}
محسوب می‌شود. این مدل با تمرین بر روی داده‌های مصنوعی، توانسته است درک و تولید متن در زبان‌های با منابع کم را به‌طور موثری بهبود ببخشد. آیا اکسپنس از داده های تولید شده توسط مدل های زبانی بزرگ دیگر به عنوان داده های تمرین استفاده میکند،  این رویکرد نوآورانه به آیا اکسپنس اجازه می‌دهد تا از منابع داده مصنوعی متنوعی استفاده کند، اما در عین حال مشکلاتی را درباره قابلیت اطمینان مدل و عدم قطعیت ذاتی در خروجی‌های مدل آیا اکپسنس و مدل هایی که بر اساس آن به وجود آمده اند(مانند مدل های گائوکرنا) به وجود می‌آورد.

یک آزمایش ساده برای نمایش عدم قطعیت موجود در هر دو مدل آیا اکسپنس و مدل گائوکرنا انجام شده است.در این آزمایش، مدل‌های آیا اکسپنس،گائوکرنا وی و
\lr{Med-Gemma}
برای پاسخگویی به سوالات کنکور علوم پزشکی پایه ایران که در شهریور
\lr{1402}
برگزار شد، به کار گرفته شدند. هر مدل پنج بار بر روی این مجموعه از سوالات اجرا شد. نتایج نشان داد که تنها
\lr{Med-Gemma}
به طور مداوم یک گزینه را برای پرسش مورد نظر در هر پنج اجرا انتخاب میکند، در حالی که مدل‌های دیگر تنوع بسیار بالایی در پاسخ از خود نشان دادند و تنها برای برخی سوالات گزینه یکسانی را در هر پنج اجرا برگزیدند. این یافته‌ها در جدول
\ref{tab:certainty_comparison}
خلاصه شده‌اند. این آزمایش، سطح عدم قطعیت مرتبط با مدل‌ آیا اکسپنس را برجسته می‌کند و چالش‌هایی را که این مدل‌ در ارائه خروجی‌های قابل اعتماد و سازگار با آن مواجه است را آشکار می‌سازد. بنابراین پرداختن به این موضوع نیز میتواند موضوع یکی از تحقیقات آینده باشد.


\begin{table}[ht]
		\centering
		\begin{tabular}{|l|c|c|c|}  
			\hline
			\textbf{} & \textbf{gaokerena-V} & \textbf{aya-expanse} & \textbf{medgemma} \\ \hline
			\lr{Accuracy}  & \textbf{\lr{38.69}} & \lr{34.52}  & \lr{37.5} \\ \hline
			\lr{Number of}& &  & \\ 
			\lr{Single Option Taken}& \lr{10}  & \lr{44}  & \textbf{\lr{168}} \\ \hline
                         \lr{Number of}& &  & \\ 
			\lr{All Option Taken}& \lr{10}  & \lr{7}  & \textbf{\lr{0}} \\ \hline
			\lr{Entropy}& \lr{1.62} & \lr{1.21}  & \textbf{\lr{0}}   \\ \hline
                            \lr{Prompt}& \lr{COT}  &  \lr{COT} &  \lr{COT} \\ \hline
                            \lr{Number of}& &  & \\ 
                            \lr{Parameters}& \lr{8 billion} & \lr{8 billion}  & \lr{4 billion}  \\ \hline
		\end{tabular}
		\caption{مقایسه اطمینان مدل ها روی مجموعه داده سوالات کنکور علوم پایه پزشکی  شهریور
		\lr{1402}
		}
		\label{tab:certainty_comparison}
	\end{table}





\begin{thebibliography}{00}
        \bibitem{b1}
        \lr{Momani, Ahmad. "Implications of Artificial Intelligence on Health Data Privacy and Confidentiality." arXiv preprint arXiv:2501.01639 (2025).‏}
        \bibitem{b2}
        \lr{Liu, Fenglin, et al. "Application of large language models in medicine." Nature Reviews Bioengineering 3.6 (2025): 445-464.‏}
        \bibitem{b3}
        \lr{Bean, Andrew M., et al. "Reliability of LLMs as medical assistants for the general public: a randomized preregistered study." Nature Medicine (2026): 1-7.‏}
        \bibitem{b4}
        \lr{Wang, Wenxuan, et al. "Medical reasoning in the era of LLMs: a systematic review of enhancement techniques and applications." arXiv preprint arXiv:2508.00669 (2025).‏}
	\bibitem{b5}
	\lr{Brown, Peter F., et al. "A statistical approach to machine translation." (1990): 79-85.}
	\bibitem{b6}
	\lr{Williams, Ronald J., and David Zipser. "A learning algorithm for continually running fully recurrent neural networks." Neural computation 1.2 (1989): 270-280.}
	\bibitem{b7}
	\lr{Hochreiter, Sepp, and Jürgen Schmidhuber. "Long short-term memory." Neural computation 9.8 (1997): 1735-1780.}
	\bibitem{b8} 
	\lr{Vaswani, Ashish, et al. "Attention is all you need." Advances in neural information processing systems 30 (2017).}
	\bibitem{b9}
	\lr{Koroteev, Mikhail V. "BERT: a review of applications in natural language processing and understanding." arXiv preprint arXiv:2103.11943 (2021).}
	\bibitem{b10}
	\lr{Devlin, Jacob, et al. "Bert: Pre-training of deep bidirectional transformers for language understanding." Proceedings of the 2019 conference of the North American chapter of the association for computational linguistics: human language technologies, volume 1 (long and short papers). 2019.}
	\bibitem{b11}
	\lr{Yenduri, Gokul, et al. "Generative pre-trained transformer: A comprehensive review on enabling technologies, potential applications, emerging challenges, and future directions." arXiv preprint arXiv:2305.10435 (2023).}
	\bibitem{b12}
	\lr{Raffel, Colin, et al. "Exploring the limits of transfer learning with a unified text-to-text transformer." Journal of machine learning research 21.140 (2020): 1-67.}
	\bibitem{b13}
	\lr{Luo, Man, et al. "Choose your QA model wisely: A systematic study of generative and extractive readers for question answering." arXiv preprint arXiv:2203.07522 (2022).}
	\bibitem{b14}
	\lr{Hendrycks, Dan, et al. "Measuring massive multitask language understanding." arXiv preprint arXiv:2009.03300 (2020).}
	\bibitem{b15} 
	\lr{Liu, Yang, et al. "G-eval: NLG evaluation using gpt-4 with better human alignment." arXiv preprint arXiv:2303.16634 (2023).}
	\bibitem{b16}
	\lr{Manes, Itay, et al. "K-qa: A real-world medical q\&a benchmark." arXiv preprint arXiv:2401.14493 (2024).}
	\bibitem{b17}
	\lr{Zhang, Tianyi, et al. "Bertscore: Evaluating text generation with bert." arXiv preprint arXiv:1904.09675 (2019).‏}
	\bibitem{b18}
	\lr{Christiano, Paul F., et al. "Deep reinforcement learning from human preferences." Advances in neural information processing systems 30 (2017).‏}
	\bibitem{b19}
	\lr{Lee, Harrison, et al. "Rlaif vs. rlhf: Scaling reinforcement learning from human feedback with ai feedback." arXiv preprint arXiv:2309.00267 (2023).‏}
	\bibitem{b20}
	\lr{Schulman, John, et al. "Proximal policy optimization algorithms." arXiv preprint arXiv:1707.06347 (2017).‏}
	\bibitem{b21}
	\lr{Shao, Zhihong, et al. "Deepseekmath: Pushing the limits of mathematical reasoning in open language models." arXiv preprint arXiv:2402.03300 (2024).‏}
	\bibitem{b22}
	\lr{Rafailov, Rafael, et al. "Direct preference optimization: Your language model is secretly a reward model." Advances in neural information processing systems 36 (2023): 53728-53741.‏}
	\bibitem{b23}
	\lr{Ji, Kaixuan, et al. "Enhancing multi-step reasoning abilities of language models through direct q-function optimization." arXiv preprint arXiv:2410.09302 (2024).‏}
	\bibitem{b24}
	\lr{Haarnoja, Tuomas, et al. "Soft actor-critic: Off-policy maximum entropy deep reinforcement learning with a stochastic actor." International conference on machine learning. Pmlr, 2018.‏}
	\bibitem{b25}
	\lr{Bengio, Yoshua. "From system 1 deep learning to system 2 deep learning." Neural Information Processing Systems. 2019.‏}
	\bibitem{b26}
	\lr{Kahneman, Daniel. "Thinking, fast and slow." Farrar, Straus and Giroux (2011).‏}
	\bibitem{b27}
	\lr{Goyal, Anirudh, et al. "Recurrent independent mechanisms." arXiv preprint arXiv:1909.10893 (2019).‏}
	\bibitem{b28}
	\lr{Wang, Guan, et al. "Hierarchical Reasoning Model." arXiv preprint arXiv:2506.21734 (2025).‏‏}
	\bibitem{b29}
	\lr{Pan, Liangming, et al. "Automatically correcting large language models: Surveying the landscape of diverse self-correction strategies." arXiv preprint arXiv:2308.03188 (2023).‏}
	\bibitem{b30}
	\lr{Liang, Baoyu, Yuchen Wang, and Chao Tong. "AI Reasoning in Deep Learning Era: From Symbolic AI to Neural–Symbolic AI." Mathematics 13.11 (2025): 1707.‏}
	\bibitem{b31}
	\lr{Lightman, Hunter, et al. "Let's verify step by step." The Twelfth International Conference on Learning Representations. 2023.}
	\bibitem{b32}
	\lr{Jiang, Shuyang, et al. "MedS $^ 3$: Towards Medical Slow Thinking with Self-Evolved Soft Dual-sided Process Supervision." arXiv preprint arXiv:2501.12051 (2025).‏}
	\bibitem{b33}
	\lr{Huang, Jie, and Kevin Chen-Chuan Chang. "Towards reasoning in large language models: A survey." Findings of the association for computational linguistics: ACL 2023. 2023.‏}
	\bibitem{b34}
	\lr{Shojaee, Parshin, et al. "The illusion of thinking: Understanding the strengths and limitations of reasoning models via the lens of problem complexity." arXiv preprint arXiv:2506.06941 (2025).‏}
	\bibitem{b35}
	\lr{Wei, Jason, et al. "Emergent abilities of large language models." arXiv preprint arXiv:2206.07682 (2022).‏}
	\bibitem{b36}
	\lr{Guo, Daya, et al. "Deepseek-r1: Incentivizing reasoning capability in llms via reinforcement learning." arXiv preprint arXiv:2501.12948 (2025).‏}
	\bibitem{b37} 
	\lr{Singhal, Karan, et al. "Toward expert-level medical question answering with large language models." Nature Medicine (2025): 1-8.}
	\bibitem{b38}
	\lr{Singhal, Karan, et al. "Large language models encode clinical knowledge." Nature 620.7972 (2023): 172-180.}
	\bibitem{b39}
	\lr{Saab, Khaled, et al. "Capabilities of gemini models in medicine." arXiv preprint arXiv:2404.18416 (2024).}
	\bibitem{b40} 
	\lr{Li, Yunxiang, et al. "Chatdoctor: A medical chat model fine-tuned on a large language model meta-ai (llama) using medical domain knowledge." Cureus 15.6 (2023).}
	\bibitem{b41} 
	\lr{Touvron, Hugo, et al. "Llama: Open and efficient foundation language models." arXiv preprint arXiv:2302.13971 (2023).}
	\bibitem{b42} 
	\lr{Kim, Hyunjae, et al. "Small language models learn enhanced reasoning skills from medical textbooks." arXiv preprint arXiv:2404.00376 (2024).}
	\bibitem{b43} 
	\lr{Vishwanath, Krithik, et al. "MedMobile: A mobile-sized language model with expert-level clinical capabilities." arXiv preprint arXiv:2410.09019 (2024).}
	\bibitem{b44} 
	\lr{Abdin, Marah, et al. "Phi-3 technical report: A highly capable language model locally on your phone." arXiv preprint arXiv:2404.14219 (2024).}
	\bibitem{b45} 
	\lr{Taghizadeh, Nasrin, et al. "SINA-BERT: a pre-trained language model for analysis of medical texts in Persian." arXiv preprint arXiv:2104.07613 (2021).}
	\bibitem{b46} 
	\lr{Koroteev, Mikhail V. "BERT: a review of applications in natural language processing and understanding." arXiv preprint arXiv:2103.11943 (2021).}
	\bibitem{b47}
	\lr{Veisi, Hadi, and Hamed Fakour Shandi. "A Persian medical question answering system." International Journal on Artificial Intelligence Tools 29.06 (2020): 2050019.}
	\bibitem{b48}
	\lr{Darabi, Leila. Medical Question Answering for Persian. Master’s thesis, LIACS, Leiden University, 2024.}
	\bibitem{b49}
	\lr{Farahani, Mehrdad, et al. "Parsbert: Transformer-based model for persian language understanding." Neural Processing Letters 53 (2021): 3831-3847.}
	\bibitem{b50}
	\lr{Coulom, Rémi. "Efficient selectivity and backup operators in Monte-Carlo tree search." International conference on computers and games. Berlin, Heidelberg: Springer Berlin Heidelberg, 2006.‏}
	\bibitem{b51}
	\lr{Wu, Juncheng, et al. "Medreason: Eliciting factual medical reasoning steps in llms via knowledge graphs." arXiv preprint arXiv:2504.00993 (2025).‏}
	\bibitem{b52}
	\lr{Chen, Junying, et al. "Huatuogpt-o1, towards medical complex reasoning with llms." arXiv preprint arXiv:2412.18925 (2024).‏}
	\bibitem{b53}
	\lr{Wu, Tianhao, et al. "Thinking llms: General instruction following with thought generation." arXiv preprint arXiv:2410.10630 (2024).‏}
	\bibitem{b54}
	\lr{Ho, Namgyu, Laura Schmid, and Se-Young Yun. "Large language models are reasoning teachers." Proceedings of the 61st annual meeting of the association for computational linguistics (volume 1: long papers). 2023.}‏
	\bibitem{b55}
	\lr{García-Ferrero, Iker, et al. "Medical mT5: an open-source multilingual text-to-text LLM for the medical domain." arXiv preprint arXiv:2404.07613 (2024).}
	\bibitem{b56}
	\lr{Liu, Yang, et al. "Datasets for large language models: A comprehensive survey." arXiv preprint arXiv:2402.18041 (2024).}
	\bibitem{b57}
	\lr{Yang, Songhua, et al. "Zhongjing: Enhancing the chinese medical capabilities of large language model through expert feedback and real-world multi-turn dialogue." Proceedings of the AAAI conference on artificial intelligence. Vol. 38. No. 17. 2024.}
	\bibitem{b58}
	\lr{Bao, Zhijie, et al. "Disc-medllm: Bridging general large language models and real-world medical consultation." arXiv preprint arXiv:2308.14346 (2023).}
	\bibitem{b59}
	\lr{Zhang, Hongbo, et al. "Huatuogpt, towards taming language model to be a doctor." arXiv preprint arXiv:2305.15075 (2023).}
	\bibitem{b60}
	\lr{Wang, Xidong, et al. "Huatuo-26M, a Large-scale Chinese Medical QA Dataset." Findings of the Association for Computational Linguistics: NAACL 2025. 2025.}
	\bibitem{b61}
	\lr{Zeng, Guangtao, et al. "MedDialog: Large-scale medical dialogue datasets." Proceedings of the 2020 conference on empirical methods in natural language processing (EMNLP). 2020.}
	\bibitem{b62}
	\lr{Han, Tianyu, et al. "MedAlpaca--an open-source collection of medical conversational AI models and training data." arXiv preprint arXiv:2304.08247 (2023).}
	\bibitem{b63}
	\lr{Bi, Xiao, et al. "Deepseek llm: Scaling open-source language models with longtermism." arXiv preprint arXiv:2401.02954 (2024).‏}
	\bibitem{b64}
	\lr{Pal, Ankit, Logesh Kumar Umapathi, and Malaikannan Sankarasubbu. "Medmcqa: A large-scale multi-subject multi-choice dataset for medical domain question answering." Conference on health, inference, and learning. PMLR, 2022.‏}
	\bibitem{b65}
	\lr{C-LARA-Instance, Manny Rayner ChatGPT. "How Woke is Grok? Empirical Evidence that xAI’s Grok Aligns Closely with Other Frontier Models." (2025).‏}
	\bibitem{b66}
	\lr{Achiam, Josh, et al. "Gpt-4 technical report." arXiv preprint arXiv:2303.08774 (2023).‏}
	\bibitem{b67}
	\lr{Jamali, Naghmeh, et al. "PerMedCQA: Benchmarking Large Language Models on Medical Consumer Question Answering in Persian Language." arXiv preprint arXiv:2505.18331 (2025).‏}
	\bibitem{b68}
	\lr{Kalahroodi, Mohammad Javad Ranjbar, et al. "PersianMedQA: Language-Centric Evaluation of LLMs in the Persian Medical Domain." arXiv preprint arXiv:2506.00250 (2025).‏}
	\bibitem{b69}
	\lr{Khoramfar, Ali, Mohammad Javad Dousti, and Heshaam Faili. "PerMed-MM: A Multimodal, Multi-Specialty Persian Medical Benchmark for Evaluating Vision Language Models." Proceedings of the 14th International Joint Conference on Natural Language Processing and the 4th Conference of the Asia-Pacific Chapter of the Association for Computational Linguistics. 2025.‏}
	\bibitem{b70}
	\lr{Yang, An, et al. “Qwen2 Technical Report.” arXiv Preprint arXiv:2407.10671, 2024.}
	\bibitem{b71} 
	\lr{Dang, John, et al. "Aya expanse: Combining research breakthroughs for a new multilingual frontier." arXiv preprint arXiv:2412.04261 (2024).}
	\bibitem{b72} 
	\lr{Team, Gemma, et al. "Gemma 2: Improving open language models at a practical size, 2024." URL https://arxiv. org/abs/2408.00118 1.3 (2024).}
	\bibitem{b73} 
	\lr{Rostami, Pedram, Ali Salemi, and Mohammad Javad Dousti. "Persianmind: A cross-lingual persian-english large language model." arXiv preprint arXiv:2401.06466 (2024).}
	\bibitem{b74}
	\lr{Dash, Saurabh, et al. "Aya Vision: Advancing the Frontier of Multilingual Multimodality." arXiv preprint arXiv:2505.08751 (2025).}
	\bibitem{b75}
	\lr{Shumailov, Ilia, et al. "AI models collapse when trained on recursively generated data." Nature 631.8022 (2024): 755-759.}
	\bibitem{b76}
	\lr{Odumakinde, Ayomide, et al. "Multilingual arbitrage: Optimizing data pools to accelerate multilingual progress, 2024." URL https://arxiv.org/abs/2408.14960.}
	\bibitem{b77}
	\lr{Hu, Edward J., et al. "Lora: Low-rank adaptation of large language models." ICLR 1.2 (2022): 3.}
	\bibitem{b78}
	\lr{Dao, Tri. "Flashattention-2: Faster attention with better parallelism and work partitioning." arXiv preprint arXiv:2307.08691 (2023).}
	\bibitem{b79}
	\lr{Lacoste, Alexandre, et al. "Quantifying the carbon emissions of machine learning." arXiv preprint arXiv:1910.09700 (2019).}
	\bibitem{b80}
	\lr{Hurst, Aaron, et al. "Gpt-4o system card." arXiv preprint arXiv:2410.21276 (2024).}
	\bibitem{b81}
	\lr{Khashabi, Daniel, et al. "Parsinlu: a suite of language understanding challenges for persian." Transactions of the Association for Computational Linguistics 9 (2021): 1147-1162.}
	\bibitem{b82}
	\lr{Kashefi, Omid. "MIZAN: a large persian-english parallel corpus." arXiv preprint arXiv:1801.02107 (2018).}
	\bibitem{b83}
	\lr{Wei, Jason, et al. "Chain-of-thought prompting elicits reasoning in large language models." Advances in neural information processing systems 35 (2022): 24824-24837.‏‏}
	\bibitem{b84}
	\lr{Brown, Bradley, et al. "Large language monkeys: Scaling inference compute with repeated sampling." arXiv preprint arXiv:2407.21787 (2024).‏}
	\bibitem{b85}
	\lr{Zhang, Mozhi, et al. "Calibrating the confidence of large language models by eliciting fidelity." arXiv preprint arXiv:2404.02655 (2024).‏}
\end{thebibliography}
